% Iteration 1 change manifest as two structured cards: chg-1 edits the system
% prompt, chg-2 edits the shell tool. Both come from the same evolve round and
% illustrate that a single round can mix harness components of different
% controllability levels. Palette and equal-height layout follow case-trajectory.tex.
\begin{figure}[t]
    \centering
    \scriptsize
    \begin{tcbraster}[raster columns=2,
                      raster equal height,
                      raster column skip=4pt,
                      raster row skip=0pt,
                      raster left skip=0pt,
                      raster right skip=0pt,
                      boxrule=0.6pt,
                      left=4pt,right=4pt,top=3pt,bottom=3pt,
                      coltitle=white]
        \begin{tcolorbox}[colback=aheSky!18,colframe=aheNavy,colbacktitle=aheNavy,
                          title={\scriptsize\textbf{chg-1，迭代 1，commit \texttt{c0b8a05}，层级：提示词}}]
            \textbf{\textcolor{aheNavy}{涉及文件}} \\
            \rule{\linewidth}{0.3pt}\\[1pt]
            $\bullet$ \texttt{workspace/systemprompt.md} \\[3pt]
            \rule{\linewidth}{0.3pt}\\
            \textbf{\textcolor{aheNavy}{改动内容}} \\
            \rule{\linewidth}{0.3pt}\\[1pt]
            追加了一套契约优先的工作流，由八条编号规则组成，涵盖验收契约抽取、评测器复刻、最小编辑语义、候选打分、泛化、时间预算、终态就绪判定以及一条停止规则。其中不含 SQLite、WAL 或任何任务特定的关键词。 \\[3pt]
            \rule{\linewidth}{0.3pt}\\
            \textbf{\textcolor{aheNavy}{修复的失败模式}} \\
            \rule{\linewidth}{0.3pt}\\[1pt]
            智能体依据自创的代理检查（例如行数或文件是否存在）就提交，而不是复现评测器字面上的断言。 \\[3pt]
            \rule{\linewidth}{0.3pt}\\
            \textbf{\textcolor{aheNavy}{预测修复的任务}} \\
            \rule{\linewidth}{0.3pt}\\[1pt]
            14 个任务。例如：\texttt{configure-git-webserver}、\texttt{query-optimize}、\texttt{mteb-retrieve}、\texttt{train-fasttext}。
        \end{tcolorbox}
        \begin{tcolorbox}[colback=aheTeal!12,colframe=aheBlue,colbacktitle=aheBlue,
                          title={\scriptsize\textbf{chg-2，迭代 1，commit \texttt{169c34c}，层级：工具}}]
            \textbf{\textcolor{aheBlue}{涉及文件}} \\
            \rule{\linewidth}{0.3pt}\\[1pt]
            $\bullet$ \texttt{tool\_descriptions/run\_shell\_command.tool.yaml} \\
            $\bullet$ \texttt{tools/shell\_tools/run\_shell\_command.py} \\[3pt]
            \rule{\linewidth}{0.3pt}\\
            \textbf{\textcolor{aheBlue}{改动内容}} \\
            \rule{\linewidth}{0.3pt}\\[1pt]
            在 shell 工具上暴露了逐次调用的 \texttt{timeout\_ms}，补充了后台执行的指引，并在超时的 shell 输出后追加一条超时恢复提示，使智能体能够改用短探测加后台任务的方式，而不是干等默认的 5 分钟。 \\[3pt]
            \rule{\linewidth}{0.3pt}\\
            \textbf{\textcolor{aheBlue}{修复的失败模式}} \\
            \rule{\linewidth}{0.3pt}\\[1pt]
            智能体把 rollout 预算耗在长时间的前台安装和 sleep 轮询循环上，反复触发默认的 5 分钟超时。 \\[3pt]
            \rule{\linewidth}{0.3pt}\\
            \textbf{\textcolor{aheBlue}{预测修复的任务}} \\
            \rule{\linewidth}{0.3pt}\\[1pt]
            8 个任务。例如：\texttt{compile-compcert}、\texttt{regex-chess}、\texttt{adaptive-rejection-sampler}。
        \end{tcolorbox}
    \end{tcbraster}
    \caption{迭代 1 中写下的两条变更清单条目，一条修改系统提示词，一条修改 shell 工具。二者出现在 Evolve Agent 产出的同一份 \texttt{change\_manifest.json} 中，随后作为具有约束力的契约进入下一轮的第 3 阶段；若其预测的修复没有兑现，归因检查就会将其回滚。}
    \label{fig:manifest-prompt-tool}
\end{figure}
